%% example.tex
%% Copyright 2026 Hannah Kirkland
%
% This work may be distributed and/or modified under the
% conditions of the LaTeX Project Public License, either version 1.3
% of this license or (at your option) any later version.
% The latest version of this license is in
%   https://www.latex-project.org/lppl.txt
% and version 1.3c or later is part of all distributions of LaTeX
% version 2008 or later.
%
% This work has the LPPL maintenance status `maintained'.
% 
% The Current Maintainer of this work is Hannah Kirkland.
%
% This work consists of the files plotweaver.cls, plotweaver-boxes.sty,
% plotweaver-character-sheet.sty, plotweaver-colors.sty,
% plotweaver-icons.sty, plotweaver-timeline.sty, character-sheet.tex, 
% character-sheet.pdf, example.tex, example.pdf, plotweaver.tex,
% plotweaver.pdf, nomodan.png, and README.md.

\documentclass[article,stormlight,timeline]{plotweaver}
\usepackage[colorlinks=true]{hyperref}
\usepackage{enumitem}
\usepackage{graphicx}
\usepackage[base]{babel}
\usetikzlibrary{chains, backgrounds}
\pgfdeclarelayer{foreground}
\pgfdeclarelayer{lines}
\pgfsetlayers{background,lines,main,foreground}
\usepackage{wrapfig}
\setlength{\intextsep}{0.25em}
\usepackage{lipsum}
  
\title{Plotweaver \LaTeX~Document Class}
\author{}
\date{}

\begin{document}
\maketitle

This file serves as a tutorial and example for how to use the \texttt{plotweaver} document class to typeset homebrew content in the style of Plotweaver RPG and Cosmere RPG games.

\section{Introduction}
The \texttt{plotweaver} document class is a free and open source resource which uses \LaTeX~to emulate the style of the Plotweaver RPG and Cosmere RPG rule books. This allows you to focus on writing and creating content without the hassle of formatting each document. To use this document class, you will need to have \LaTeX~installed, which can be downloaded at the \href{https://www.latex-project.org/}{official \LaTeX~project website}. If you're not familiar with \LaTeX, a full installation is recommended. Another non-local option is to upload the class and package files to an online \LaTeX~editor such as Overleaf. Please note that online editors may or maynot be free. Regardless of installation choice, you will need to use Lua\LaTeX~as the compiler.

The philosophy of this document class and its accompanying packages is to be as simple to use as possible while remaining as modular and flexible as possible. It is also still under active development, particularly with respect to the talent tree boxes. Contributions are welcome on this project's \href{https://github.com/teawizardry/plotweaver}{GitHub repository}. 

This work is licensed under the \href{https://latex-project.org/lppl/lppl-1-3c/}{LaTeX Project Public License (LPPL) version 1.3c}. This is unofficial fan content, created and shared for non-commercial use. It has not been reviewed by Dragonsteel Entertainment, LLC or Brotherwise Games, LLC.

\section{Example Usage}
For full documentation, please see \texttt{plotweaver.pdf}. For an example of how to use this document class to generate character sheets, please see \texttt{character-sheet.tex}. 

\subsection{Feature Boxes}
Three feature boxes are specified using \texttt{tcolorbox}. They are \texttt{featureBox}, \texttt{leftBox}, \texttt{rightBox}, and \texttt{calloutBox}. You can also use the command \texttt{statBox} to nicely align a left and right box. These boxes support upper and lower content, but you don't have to specify both.

\subsection{Talent Paths}
TikZ is used to make the talent paths. You can add a photo as the background, if desired. Check the \texttt{example.tex} for a sample talent path.

\subsection{Icons}
Icons for activations and plot die are specified with \texttt{forever}, \texttt{freeAction}, \texttt{action}, \texttt{doubleAction}, \texttt{tripleAction}, \texttt{reaction}, \texttt{opportunity}, and \texttt{complication}. They look like this:
\begin{center}
    \forever \freeAction \action \doubleAction \tripleAction \reaction \opportunity \complication
\end{center}
The complication commands can take a number or be left blank. Numbers look like this: \complication[2]. You can also make larger complication symbols with \texttt{complicationDice}:
\begin{center}
    \complicationDice\space \complicationDice[2] \complicationDice[4]
\end{center}

\subsection{Tables}
Tables can be created using \texttt{plotable}. This uses \texttt{nicematrix} under the hood. Like usual, you can refer to Table~\ref{tab:example} with a label.
\begin{table}
    \caption{Example Table}
    \begin{plotable}{llll}
        Heading 1 & Heading 2 & Heading 3 & Heading 4 \\
        Lipsum & Lorem ipsum dolor sit amet, consectetur adipiscing elit. & 1234 & 5678 \\
        Lipsum & Lorem ipsum dolor sit amet, consectetur adipiscing elit. & 1234 & 5678 \\
        Lipsum & Lorem ipsum dolor sit amet, consectetur adipiscing elit. & 1234 & 5678 \\
        Lipsum & Lorem ipsum dolor sit amet, consectetur adipiscing elit. & 1234 & 5678 \\
    \end{plotable}
    \lowercaption{This is a caption that is placed below the table}
    \label{tab:example}
\end{table}

\subsection{Color Schemes}
You can switch color schemes by passing an option to the document class. The default is \texttt{stormlight}. You can also use \texttt{mistborn} or \texttt{ariandor}. You can also define your own color schemes.

\subsection{Character Sheets}
You can make a character sheet using the \texttt{MakeCharacterSheet} command if \texttt{character-sheet} is passed into the document class options. It will calculate derived attributes for you, and you can set it to use the terminology for Stormlight Archive or Mistborn instead of the generic Plotweaver terms. See \texttt{character-sheet.tex} for an example of how to use this command. In the future, I intend to make the sheet more customizable.

\begin{center}
\noindent
\begin{tcolorbox}[featureBox, title={Feature Box}]
    \lipsum[1][1-7]
    \tcblower
    \tcbsubtitle{Sub Feature}
    \lipsum[1][1-3]
\end{tcolorbox}
\statBox{Stat One}{
    \normalsize These boxes can have different content...
}%
{Stat Two}{
    \normalsize ...and still end up the same height! Even when one is multiple lines.
}
\begin{tcolorbox}[calloutBox, title={Callout Box}]
    This is a callout box. It will be as wide as the text (i.e., \texttt{\textbackslash linewidth}) unless you pass a width to the \texttt{tcolorbox} environment.

\end{tcolorbox}
\vspace{\baselineskip}
\end{center}

You can also have a callout box that will float next to the text. To do this, you need to use the \texttt{wrapfig} package. Set the width in the \texttt{wrapfigure} environment and then let the \texttt{tcolorbox} fill the total width.

\begin{wrapfigure}{r}{0.45\textwidth}
    \begin{tcolorbox}[calloutBox, title={Callout Box}, nobeforeafter]
        \lipsum[1][1-4]
    \end{tcolorbox}
\end{wrapfigure}

\lipsum[1]



\clearpage

\begin{figure*}
\begin{center}
\begin{tikzpicture}
 [node distance=1em and -4em,
 start chain=going below, overlay, remember picture, background rectangle/.style={fill=cscolor}, show background rectangle]

    \node (kt) [fill=none] at (0,0) {
        \keyTalent{Key Talent}{
        This is how you make a Key Talent.
        }
    };

    \node (ta1) [talent, below left = of kt] {
        \talent{\forever Forever}{
            This is how you make a Talent box.
        }
    };

    \node (ta2) [talent, below = of ta1] {
        \talent{\freeAction Free Action}[Prerequisite: Statistic +2]{
            You can also specify a prerequisite for a Talent.
        }
    };

    \node (ta3) [talent, below = of ta2] {
        \talent{\action Action}{
            You can specify complications with a number \complication[2] or with no numbers \complication. You can also make dice style complications like this \complicationDice[4]
        }
    };

    \node (tb1) [talent, below right = of kt] {
        \talent{\doubleAction Double Action}{
            You can also specify when \opportunity opportunities happen. 
        }
    };

    \node (tb2) [talent, below = of tb1] {
        \talent{\tripleAction Triple Action}{
        To make a Talent Tree, first specify all of the Talent nodes.
        }
    };

    \node (tb3) [talent, below = of tb2] {
        \talent{\reaction Reaction}{
        Then, you can use lines to connect the nodes.
        }
    };

% Connectors
\begin{pgfonlayer}{lines}
    \draw [connector] (kt.west) ++(0.15,0) -| (ta1);
    \draw [connector] (ta1) -- (ta2);
    \draw [connector] (ta2) -- (ta3);
    \draw [connector] (kt.east) ++(-0.25,0) -| (tb1);
    \draw [connector] (tb1) -- (tb2);
    \draw [connector] (tb2) -- (tb3);
    \draw [connector] (tb2.west) -| ++(-1.25,0) |- (ta3.east);
\end{pgfonlayer}
\end{tikzpicture}
\end{center}
\end{figure*}

\null\newpage

\subsection{Timelines}
Timelines can be created using the \texttt{timeline} environment if \texttt{timeline} is passed into the document class options. You can pass in arguments to determine the spacing between major and minor ticks with \texttt{major tick pattern} and adjust the overall scale with \texttt{scale}. Within the timeline environment, the start and end years are specified with the \texttt{era} command. Specific events are added aftwards with the \texttt{event} command. This allows for multiple eras to be chained together in a single timeline. 
\begin{figure}
    \begin{timeline}[major tick pattern=5, scale=0.2]
        \era{start=1700, end=1760, prefix=XX}  
        \event{year=1702}{An important event happened here.}
        \event{year=1750, prefix={The Beginning of an Era --- }}{Prefixes give different colored headings.}
        \event{year=1757}{Another important event happened here.}
    \end{timeline}
\end{figure}

\end{document}